% Figure: De Morgan's law read as a reliability statement. A two-component
% parallel-redundant system succeeds when either component works (OR); by De
% Morgan it fails only when both fail (AND). The left panel shows the success
% event as a union of two circles; the right panel shows the failure event as
% the intersection of the two complements, the shaded lens outside both.
\begin{figure}[htbp]
\centering
\begin{tikzpicture}[scale=1.0]
% ---- Left panel: success = A OR B (union) ----
\begin{scope}[shift={(0,0)}]
  \draw[engblack!55] (-1.7,-1.7) rectangle (1.9,1.7);
  \node[font=\footnotesize, anchor=north west] at (-1.65,1.62) {universe};
  \fill[engblue!22] (-0.45,0) circle (0.95);
  \fill[engblue!22] (0.45,0) circle (0.95);
  \draw[engblue, thick] (-0.45,0) circle (0.95);
  \draw[engblue, thick] (0.45,0) circle (0.95);
  \node[font=\footnotesize, engblue] at (-0.95,0) {$A$};
  \node[font=\footnotesize, engblue] at (0.95,0) {$B$};
  \node[font=\small, anchor=north] at (0.1,-1.95)
    {success $=A\cup B$};
\end{scope}
% ---- Right panel: failure = (not A) AND (not B) = complement of union ----
\begin{scope}[shift={(5.2,0)}]
  \fill[engred!16] (-1.7,-1.7) rectangle (1.9,1.7);
  \fill[white] (-0.45,0) circle (0.95);
  \fill[white] (0.45,0) circle (0.95);
  \draw[engblack!55] (-1.7,-1.7) rectangle (1.9,1.7);
  \draw[engred, thick] (-0.45,0) circle (0.95);
  \draw[engred, thick] (0.45,0) circle (0.95);
  \node[font=\footnotesize, engred] at (-0.95,0) {$A$};
  \node[font=\footnotesize, engred] at (0.95,0) {$B$};
  \node[font=\footnotesize, anchor=north east] at (1.85,1.62)
    {$\overline{A}\cap\overline{B}$};
  \node[font=\small, anchor=north] at (0.1,-1.95)
    {failure $=\overline{A\cup B}$};
\end{scope}
\end{tikzpicture}
\caption[De Morgan's law as a redundancy statement]{De Morgan's law read as
the reliability of a two-component parallel-redundant system. The event $A$
is ``component one works'' and $B$ is ``component two works.'' The left panel
shades the success event, the union $A\cup B$: the system works when either
component works. The right panel shades its complement, the failure event
$\overline{A\cup B}$, which De Morgan's law rewrites as $\overline{A}\cap
\overline{B}$: the system fails only in the region outside both circles,
where component one fails and component two fails together. Reading the
picture as probabilities and assuming independence, the failure probability
is the product of the two individual failure probabilities, which is the
parallel-redundancy calculation of Volume~I. The operation flips from OR to
AND exactly when the complement crosses it, which is the content of the law
and the reason a reliability statement about success and a reliability
statement about failure are the same statement seen from two sides.}
\label{fig:vol02:ch01:de-morgan-redundancy}
\end{figure}
